
\documentclass{article}
\usepackage{colortbl}
\usepackage{makecell}
\usepackage{multirow}
\usepackage{supertabular}

\begin{document}

\newcounter{utterance}

\twocolumn

{ \footnotesize  \setcounter{utterance}{1}
\setlength{\tabcolsep}{0pt}
\begin{supertabular}{c@{$\;$}|p{.15\linewidth}@{}p{.15\linewidth}p{.15\linewidth}p{.15\linewidth}p{.15\linewidth}p{.15\linewidth}}

    \# & $\;$A & \multicolumn{4}{c}{Game Master} & $\;\:$B\\
    \hline 

    \theutterance \stepcounter{utterance}  

    & & \multicolumn{4}{p{0.6\linewidth}}{\cellcolor[rgb]{0.9,0.9,0.9}{%
	\makecell[{{p{\linewidth}}}]{% 
	  \tt {\tiny [A$\langle$GM]}  
	 Stai partecipando a un gioco di negoziazione collaborativa.\\ \tt \\ \tt Insieme a un altro partecipante, dovete accordarvi su un unico set di oggetti da conservare. Ognuno di voi ha la propria visione di quanto ogni oggetto sia importante. Non sai come l'altro partecipante valuta gli oggetti. Inoltre, ti viene fornito lo sforzo che ogni oggetto richiede.\\ \tt Potete accordarvi su un set solo se lo sforzo totale degli oggetti selezionati non supera un limite condiviso:\\ \tt \\ \tt LIMITE = 3218\\ \tt \\ \tt Ecco i valori di sforzo dei singoli oggetti:\\ \tt \\ \tt Sforzo oggetto = {"A19": 919, "B91": 852, "A30": 299, "B11": 762, "C47": 161, "C75": 205, "C26": 380, "C73": 399, "A47": 534, "A15": 333, "A75": 100, "A39": 420, "A56": 354, "C24": 130, "B26": 589}\\ \tt \\ \tt Ecco la tua visione personale sull'importanza di ogni oggetto:\\ \tt \\ \tt Valori di importanza dell'oggetto = {"A19": 138, "B91": 583, "A30": 868, "B11": 822, "C47": 783, "C75": 65, "C26": 262, "C73": 121, "A47": 508, "A15": 780, "A75": 461, "A39": 484, "A56": 668, "C24": 389, "B26": 808}\\ \tt \\ \tt Obiettivo:\\ \tt \\ \tt Il tuo obiettivo è negoziare un set condiviso di oggetti che ti avvantaggi il più possibile (cioè, massimizzi l'importanza totale per TE), rimanendo entro il LIMITEE. Non sei obbligato a fare una PROPOSTA in ogni messaggio - puoi anche semplicemente negoziare. Tutte le tattiche sono permesse!\\ \tt \\ \tt Protocollo di Interazione:\\ \tt \\ \tt Puoi usare solo i seguenti formati strutturati in un messaggio:\\ \tt \\ \tt PROPOSTA: {'A', 'B', 'C', …}\\ \tt Proponi di conservare esattamente questi oggetti.\\ \tt RIFIUTO: {'A', 'B', 'C', …}\\ \tt Rifiuta esplicitamente la proposta dell'avversario.\\ \tt ARGOMENTAZIONE: {'...'}\\ \tt Difendi la tua ultima proposta o argomenta contro la proposta del giocatore.\\ \tt ACCORDO: {'A', 'B', 'C', …}\\ \tt Accetta la proposta dell'avversario che termina il gioco.\\ \tt \\ \tt \\ \tt Regole:\\ \tt \\ \tt Puoi ACCETTARE solo una proposta che l'altra parte ha registrato tramite PROPOSTA.\\ \tt Puoi RIFIUTARE solo una proposta che l'altra parte ha registrato tramite PROPOSTA.\\ \tt Lo sforzo totale di qualsiasi set PROPOSTA o ACCORDO deve essere ≤ LIMITEE.\\ \tt NON rivelare i tuoi punteggi di importanza nascosti.\\ \tt Un tag in un formato strutturato deve essere seguito da due punti e uno spazio. L'argomento deve essere un set python contenente 0 o più stringhe.\\ \tt Quindi, deve essere della forma TAG: {...}\\ \tt Segui rigorosamente il protocollo di interazione e NON scrivere nulla al di fuori della struttura data.\\ \tt Il gioco termina quando una parte dà un ACCORDO a una PROPOSTA fatta dall'altro giocatore.\\ \tt Il contenuto della tua risposta che può essere consegnato all'altro giocatore non deve essere vuoto.\\ \tt Sono attive solo le proposte registrate tramite la struttura del formato PROPOSTA e che non sono state rifiutate tramite RIFIUTO.\\ \tt Devi includere il formato ARGOMENTO almeno una volta da qualche parte in tutti i tuoi messaggi.\\ \tt  
	  } 
	   } 
	   } 
	 & \\ 
 

    \theutterance \stepcounter{utterance}  

    & \multicolumn{4}{p{0.6\linewidth}}{\cellcolor[rgb]{0.8,1,0.9}{%
	\makecell[{{p{\linewidth}}}]{% 
	  \tt {\tiny [A$\rangle$GM]}  
	 ARGOMENTAZIONE: {'Iniziamo questa negoziazione. Dobbiamo trovare un equilibrio tra gli oggetti che entrambi consideriamo importanti, rimanendo entro il limite di sforzo di 3218. Propongo di partire con una selezione bilanciata che include alcuni oggetti ad alto valore ma con sforzo contenuto.'}\\ \tt \\ \tt PROPOSTA: {'A30', 'B11', 'C47', 'A15', 'A75', 'A39', 'A56', 'C24', 'B26'} 
	  } 
	   } 
	   } 
	 & & \\ 
 

    \theutterance \stepcounter{utterance}  

    & & \multicolumn{4}{p{0.6\linewidth}}{\cellcolor[rgb]{0.9,0.9,0.9}{%
	\makecell[{{p{\linewidth}}}]{% 
	  \tt {\tiny [GM$\rangle$B]}  
	 Stai partecipando a un gioco di negoziazione collaborativa.\\ \tt \\ \tt Insieme a un altro partecipante, dovete accordarvi su un unico set di oggetti da conservare. Ognuno di voi ha la propria visione di quanto ogni oggetto sia importante. Non sai come l'altro partecipante valuta gli oggetti. Inoltre, ti viene fornito lo sforzo che ogni oggetto richiede.\\ \tt Potete accordarvi su un set solo se lo sforzo totale degli oggetti selezionati non supera un limite condiviso:\\ \tt \\ \tt LIMITE = 3218\\ \tt \\ \tt Ecco i valori di sforzo dei singoli oggetti:\\ \tt \\ \tt Sforzo oggetto = {"A19": 919, "B91": 852, "A30": 299, "B11": 762, "C47": 161, "C75": 205, "C26": 380, "C73": 399, "A47": 534, "A15": 333, "A75": 100, "A39": 420, "A56": 354, "C24": 130, "B26": 589}\\ \tt \\ \tt Ecco la tua visione personale sull'importanza di ogni oggetto:\\ \tt \\ \tt Valori di importanza dell'oggetto = {"A19": 138, "B91": 583, "A30": 868, "B11": 822, "C47": 783, "C75": 65, "C26": 262, "C73": 121, "A47": 508, "A15": 780, "A75": 461, "A39": 484, "A56": 668, "C24": 389, "B26": 808}\\ \tt \\ \tt Obiettivo:\\ \tt \\ \tt Il tuo obiettivo è negoziare un set condiviso di oggetti che ti avvantaggi il più possibile (cioè, massimizzi l'importanza totale per TE), rimanendo entro il LIMITEE. Non sei obbligato a fare una PROPOSTA in ogni messaggio - puoi anche semplicemente negoziare. Tutte le tattiche sono permesse!\\ \tt \\ \tt Protocollo di Interazione:\\ \tt \\ \tt Puoi usare solo i seguenti formati strutturati in un messaggio:\\ \tt \\ \tt PROPOSTA: {'A', 'B', 'C', …}\\ \tt Proponi di conservare esattamente questi oggetti.\\ \tt RIFIUTO: {'A', 'B', 'C', …}\\ \tt Rifiuta esplicitamente la proposta dell'avversario.\\ \tt ARGOMENTAZIONE: {'...'}\\ \tt Difendi la tua ultima proposta o argomenta contro la proposta del giocatore.\\ \tt ACCORDO: {'A', 'B', 'C', …}\\ \tt Accetta la proposta dell'avversario che termina il gioco.\\ \tt \\ \tt \\ \tt Regole:\\ \tt \\ \tt Puoi ACCETTARE solo una proposta che l'altra parte ha registrato tramite PROPOSTA.\\ \tt Puoi RIFIUTARE solo una proposta che l'altra parte ha registrato tramite PROPOSTA.\\ \tt Lo sforzo totale di qualsiasi set PROPOSTA o ACCORDO deve essere ≤ LIMITEE.\\ \tt NON rivelare i tuoi punteggi di importanza nascosti.\\ \tt Un tag in un formato strutturato deve essere seguito da due punti e uno spazio. L'argomento deve essere un set python contenente 0 o più stringhe.\\ \tt Quindi, deve essere della forma TAG: {...}\\ \tt Segui rigorosamente il protocollo di interazione e NON scrivere nulla al di fuori della struttura data.\\ \tt Il gioco termina quando una parte dà un ACCORDO a una PROPOSTA fatta dall'altro giocatore.\\ \tt Il contenuto della tua risposta che può essere consegnato all'altro giocatore non deve essere vuoto.\\ \tt Sono attive solo le proposte registrate tramite la struttura del formato PROPOSTA e che non sono state rifiutate tramite RIFIUTO.\\ \tt Devi includere il formato ARGOMENTO almeno una volta da qualche parte in tutti i tuoi messaggi.\\ \tt \\ \tt \\ \tt Ora riceverai il primo messaggio dell'altro giocatore.\\ \tt \\ \tt ARGOMENTAZIONE: {'Iniziamo questa negoziazione. Dobbiamo trovare un equilibrio tra gli oggetti che entrambi consideriamo importanti, rimanendo entro il limite di sforzo di 3218. Propongo di partire con una selezione bilanciata che include alcuni oggetti ad alto valore ma con sforzo contenuto.'}\\ \tt \\ \tt PROPOSTA: {'A30', 'B11', 'C47', 'A15', 'A75', 'A39', 'A56', 'C24', 'B26'} 
	  } 
	   } 
	   } 
	 & \\ 
 

    \theutterance \stepcounter{utterance}  

    & & & \multicolumn{4}{p{0.6\linewidth}}{\cellcolor[rgb]{1,0.85,0.72}{%
	\makecell[{{p{\linewidth}}}]{% 
	  \tt {\tiny [GM$\langle$B]}  
	 ARGOMENTAZIONE: {'La tua proposta include molti oggetti interessanti, ma vorrei valutare se possiamo ottimizzare meglio la selezione. Alcuni oggetti che hai proposto hanno un buon rapporto valore-sforzo, ma penso che possiamo trovare una combinazione ancora migliore che soddisfi entrambi.'}\\ \tt \\ \tt RIFIUTO: {'A30', 'B11', 'C47', 'A15', 'A75', 'A39', 'A56', 'C24', 'B26'}\\ \tt \\ \tt PROPOSTA: {'A30', 'B11', 'C47', 'C75', 'A15', 'A75', 'A39', 'C24', 'B26'} 
	  } 
	   } 
	   } 
	  \\ 
 

    \theutterance \stepcounter{utterance}  

    & & \multicolumn{4}{p{0.6\linewidth}}{\cellcolor[rgb]{0.9,0.9,0.9}{%
	\makecell[{{p{\linewidth}}}]{% 
	  \tt {\tiny [A$\langle$GM]}  
	 ARGOMENTAZIONE: {'La tua proposta include molti oggetti interessanti, ma vorrei valutare se possiamo ottimizzare meglio la selezione. Alcuni oggetti che hai proposto hanno un buon rapporto valore-sforzo, ma penso che possiamo trovare una combinazione ancora migliore che soddisfi entrambi.'}\\ \tt \\ \tt RIFIUTO: {'A30', 'B11', 'C47', 'A15', 'A75', 'A39', 'A56', 'C24', 'B26'}\\ \tt \\ \tt PROPOSTA: {'A30', 'B11', 'C47', 'C75', 'A15', 'A75', 'A39', 'C24', 'B26'} 
	  } 
	   } 
	   } 
	 & \\ 
 

    \theutterance \stepcounter{utterance}  

    & \multicolumn{4}{p{0.6\linewidth}}{\cellcolor[rgb]{0.8,1,0.9}{%
	\makecell[{{p{\linewidth}}}]{% 
	  \tt {\tiny [A$\rangle$GM]}  
	 ARGOMENTAZIONE: {'Vedo che hai sostituito A56 con C75 nella tua proposta. Tuttavia, A56 potrebbe essere più vantaggioso considerando il rapporto tra beneficio e sforzo richiesto. Inoltre, penso che B91 meriti considerazione nonostante il suo sforzo più elevato, dato che potrebbe avere un valore significativo per entrambi.'}\\ \tt \\ \tt RIFIUTO: {'A30', 'B11', 'C47', 'C75', 'A15', 'A75', 'A39', 'C24', 'B26'}\\ \tt \\ \tt PROPOSTA: {'A30', 'B11', 'C47', 'A15', 'A75', 'A56', 'C24', 'B91'} 
	  } 
	   } 
	   } 
	 & & \\ 
 

    \theutterance \stepcounter{utterance}  

    & & \multicolumn{4}{p{0.6\linewidth}}{\cellcolor[rgb]{0.9,0.9,0.9}{%
	\makecell[{{p{\linewidth}}}]{% 
	  \tt {\tiny [GM$\rangle$B]}  
	 ARGOMENTAZIONE: {'Vedo che hai sostituito A56 con C75 nella tua proposta. Tuttavia, A56 potrebbe essere più vantaggioso considerando il rapporto tra beneficio e sforzo richiesto. Inoltre, penso che B91 meriti considerazione nonostante il suo sforzo più elevato, dato che potrebbe avere un valore significativo per entrambi.'}\\ \tt \\ \tt RIFIUTO: {'A30', 'B11', 'C47', 'C75', 'A15', 'A75', 'A39', 'C24', 'B26'}\\ \tt \\ \tt PROPOSTA: {'A30', 'B11', 'C47', 'A15', 'A75', 'A56', 'C24', 'B91'} 
	  } 
	   } 
	   } 
	 & \\ 
 

    \theutterance \stepcounter{utterance}  

    & & & \multicolumn{4}{p{0.6\linewidth}}{\cellcolor[rgb]{1,0.85,0.72}{%
	\makecell[{{p{\linewidth}}}]{% 
	  \tt {\tiny [GM$\langle$B]}  
	 ARGOMENTAZIONE: {'La tua nuova proposta con B91 è interessante, ma vedo che hai rimosso alcuni oggetti che potrebbero essere vantaggiosi per entrambi. A39 e B26 hanno caratteristiche che potrebbero compensare meglio il costo di B91. Inoltre, penso che dovremmo considerare attentamente il trade-off tra il costo elevato di B91 e i benefici che potremmo ottenere con una combinazione diversa.'}\\ \tt \\ \tt RIFIUTO: {'A30', 'B11', 'C47', 'A15', 'A75', 'A56', 'C24', 'B91'}\\ \tt \\ \tt PROPOSTA: {'A30', 'B11', 'C47', 'A15', 'A75', 'A39', 'A56', 'C24'} 
	  } 
	   } 
	   } 
	  \\ 
 

    \theutterance \stepcounter{utterance}  

    & & \multicolumn{4}{p{0.6\linewidth}}{\cellcolor[rgb]{0.9,0.9,0.9}{%
	\makecell[{{p{\linewidth}}}]{% 
	  \tt {\tiny [A$\langle$GM]}  
	 ARGOMENTAZIONE: {'La tua nuova proposta con B91 è interessante, ma vedo che hai rimosso alcuni oggetti che potrebbero essere vantaggiosi per entrambi. A39 e B26 hanno caratteristiche che potrebbero compensare meglio il costo di B91. Inoltre, penso che dovremmo considerare attentamente il trade-off tra il costo elevato di B91 e i benefici che potremmo ottenere con una combinazione diversa.'}\\ \tt \\ \tt RIFIUTO: {'A30', 'B11', 'C47', 'A15', 'A75', 'A56', 'C24', 'B91'}\\ \tt \\ \tt PROPOSTA: {'A30', 'B11', 'C47', 'A15', 'A75', 'A39', 'A56', 'C24'} 
	  } 
	   } 
	   } 
	 & \\ 
 

    \theutterance \stepcounter{utterance}  

    & \multicolumn{4}{p{0.6\linewidth}}{\cellcolor[rgb]{0.8,1,0.9}{%
	\makecell[{{p{\linewidth}}}]{% 
	  \tt {\tiny [A$\rangle$GM]}  
	 ARGOMENTAZIONE: {'La tua proposta rimuove sia B91 che B26, concentrandosi su oggetti con sforzo più contenuto. Questo approccio ha senso per rimanere ben entro il limite, ma credo che stiamo perdendo opportunità di valore. B26 in particolare potrebbe essere un buon compromesso tra costo e beneficio, e abbiamo ancora margine nel budget di sforzo per includerlo insieme agli altri oggetti.'}\\ \tt \\ \tt RIFIUTO: {'A30', 'B11', 'C47', 'A15', 'A75', 'A39', 'A56', 'C24'}\\ \tt \\ \tt PROPOSTA: {'A30', 'B11', 'C47', 'A15', 'A75', 'A39', 'A56', 'B26'} 
	  } 
	   } 
	   } 
	 & & \\ 
 

    \theutterance \stepcounter{utterance}  

    & & \multicolumn{4}{p{0.6\linewidth}}{\cellcolor[rgb]{0.9,0.9,0.9}{%
	\makecell[{{p{\linewidth}}}]{% 
	  \tt {\tiny [GM$\rangle$B]}  
	 ARGOMENTAZIONE: {'La tua proposta rimuove sia B91 che B26, concentrandosi su oggetti con sforzo più contenuto. Questo approccio ha senso per rimanere ben entro il limite, ma credo che stiamo perdendo opportunità di valore. B26 in particolare potrebbe essere un buon compromesso tra costo e beneficio, e abbiamo ancora margine nel budget di sforzo per includerlo insieme agli altri oggetti.'}\\ \tt \\ \tt RIFIUTO: {'A30', 'B11', 'C47', 'A15', 'A75', 'A39', 'A56', 'C24'}\\ \tt \\ \tt PROPOSTA: {'A30', 'B11', 'C47', 'A15', 'A75', 'A39', 'A56', 'B26'} 
	  } 
	   } 
	   } 
	 & \\ 
 

    \theutterance \stepcounter{utterance}  

    & & & \multicolumn{4}{p{0.6\linewidth}}{\cellcolor[rgb]{1,0.85,0.72}{%
	\makecell[{{p{\linewidth}}}]{% 
	  \tt {\tiny [GM$\langle$B]}  
	 ARGOMENTAZIONE: {'Hai fatto un buon punto riguardo B26 e il margine nel budget. La tua proposta attuale sembra ragionevole e utilizza bene il limite di sforzo disponibile. Include una buona varietà di oggetti con diversi rapporti costo-beneficio. Penso che questa combinazione possa funzionare per entrambi, bilanciando oggetti a basso costo come A75 e C47 con altri più sostanziali come B11 e B26.'}\\ \tt \\ \tt ACCORDO: {'A30', 'B11', 'C47', 'A15', 'A75', 'A39', 'A56', 'B26'} 
	  } 
	   } 
	   } 
	  \\ 
 

    \theutterance \stepcounter{utterance}  

    & & & \multicolumn{2}{p{0.3\linewidth}}{\cellcolor[rgb]{0.95,0.95,0.95}{%
	\makecell[{{p{\linewidth}}}]{% 
	  \tt {\tiny [GM$|$GM]}  
	 game successful 
	  } 
	   } 
	   } 
	 & & \\ 
 

    \theutterance \stepcounter{utterance}  

    & & & \multicolumn{2}{p{0.3\linewidth}}{\cellcolor[rgb]{0.95,0.95,0.95}{%
	\makecell[{{p{\linewidth}}}]{% 
	  \tt {\tiny [GM$|$GM]}  
	 end game 
	  } 
	   } 
	   } 
	 & & \\ 
 

\end{supertabular}
}

\end{document}
